\documentclass[11pt]{amsart}
\usepackage{calc,amssymb,amsmath,ulem,stmaryrd,fullpage}
\usepackage{enumerate}
\usepackage{alltt}
\RequirePackage[dvipsnames,usenames]{color}
\let\mffont=\sf
\normalem
%\input{kmacros3.sty}
\input{xy}
\xyoption{all}
\usepackage{hyperref}
\usepackage{graphicx}
\usepackage[all,cmtip]{xy}
\usepackage{verbatim}
\usepackage[left=0.95in,top=0.95in,right=0.95in,bottom=0.95in]{geometry}
\newcommand{\tauCohomology}{T}
\newcommand{\FCohomology}{S}


\theoremstyle{theorem}
%\newtheorem*{mainthm*}{Main Theorem}

\newcommand{\note}[1]{\marginpar{\sffamily\tiny #1}}

\def\todo#1{\textcolor{Mahogany}%
{\sffamily\footnotesize\newline{\color{Mahogany}\fbox{\parbox{\textwidth-15pt}{#1}}}\newline}}

\renewcommand{\O}{\mathcal O}
\newcommand{\ie}{{\itshape i.e.} }
\newcommand{\roundup}[1]{\lceil #1 \rceil}
\newcommand{\rounddown}[1]{\lfloor #1 \rfloor}
\renewcommand{\to}[1][]{\xrightarrow{\ #1\ }}
\newcommand{\onto}[1][]{\protect{\xrightarrow{\ #1\ }\hspace{-0.8em}\rightarrow}}
\newcommand{\into}[1][]{\lhook \joinrel \xrightarrow{\ #1\ }}
%\newcommand{DBDual}[1]{{\underline{\omega}_{#1}}}
\newcommand{\DBDual}[1]{{{\uuline \omega}{}^{\mydot}_{#1}}}
\newcommand{\Tt}{{\mathfrak{T}}}
%\newcommand{\bn}{{\mathfrak{n}} }
\newcommand{\sol}[1]{ \vskip 6pt{\color{blue} {\bf Solution:} #1} \vskip 6pt}

\newcommand{\bR}{{\mathbb{R}}}
\newcommand{\va}{{\vec{a}}}
\newcommand{\vb}{{\vec{b}}}
\newcommand{\vzero}{{\vec{0}}}
\newcommand{\vi}{{\vec{i}}}
\newcommand{\vj}{{\vec{j}}}
\newcommand{\vk}{{\vec{k}}}
\newcommand{\vh}{{\vec{h}}}
\newcommand{\vr}{{\vec{r}}}
\newcommand{\vu}{{\vec{u}}}
\newcommand{\vv}{{\vec{v}}}
\newcommand{\vw}{{\vec{w}}}
\newcommand{\vx}{{\vec{x}}}
\newcommand{\vy}{{\vec{y}}}
\newcommand{\vz}{{\vec{z}}}
\newcommand{\vT}{{\vec{T}}}
\newcommand{\vN}{{\vec{N}}}
\newcommand{\vF}{{\vec{F}}}
\newcommand{\vG}{{\vec{G}}}
\newcommand{\bF}{\mathbb{F}}
\newcommand{\bP}{\mathbb{P}}
\newcommand{\bQ}{\mathbb{Q}}
\newcommand{\bZ}{\mathbb{Z}}
\newcommand{\bA}{\mathbb{A}}
\newcommand{\codim}{{\mathrm{codim}}}
\newcommand{\Map}{{\mathrm{Map}}}
\newcommand{\tld}{\widetilde}
\newcommand{\fram}{\mathfrak{m}}


\begin{document}
\title{Worksheet \#7 -- Math 6130\\Fall 2018}
\author{Due Monday, November 12th}

\maketitle
You may work in groups of up to 3.  Only one worksheet needs to be turned in per group.
\vskip 9pt
\noindent
{\bf Setup.}  Suppose that $X \subseteq \bA^n$ is an affine variety and $I \subseteq k[X]$ is an ideal (cutting out an closed subset $Z(I) = Z \subseteq X$).  Suppose $J \subseteq k[\bA^m]$ is an ideal with $Y = Z(J)$.  Finally suppose we have an injective\footnote{This isn't really necessary, but it makes life easier and the geometry simpler.} $k$-algebra map $\psi : k[\bA^m]/J \hookrightarrow k[X]/I$ that makes $k[X]/I$ into a (usually finite) $k[\bA^m]/J$-module.
\vskip 3pt
Consider the subring of $k[X]$:
\[
A := \{ f \in k[X] \;|\; \overline{f} \in \psi(k[\bA^m]/J) \subseteq k[X]/I \}.
\]
In other words, $A$ is elements of $k[X]$ whose image in $k[X]/I$ is contained the image of $k[\bA^m]/J$ in $k[X]/I$.
\vskip 9pt
\noindent
{\bf 1.}  Show that there is a commutative diagram of $k$-algebra morphisms:
\[
\xymatrix{
k[X] \ar@{->>}[r] & k[X]/I \\
A \ar@{^{(}->}[u]^{\alpha} \ar[r]_{\beta} & k[\bA^m]/J \ar@{^{(}->}[u]_{\psi}.
}
\]
\vskip 150pt
\noindent
{\bf 2.}  Show that the map $\beta$ above is surjective.
\newpage
\noindent
{\bf 3.}  Show that the fraction field of $A$ is the same as the fraction field of $k[X]$.
\vskip 3pt
\emph{Hint:} Show that every element of $I$ is in $A$.
\vskip 150pt
\noindent
{\bf 4.}  Assume $\psi$ above is finite.  Show that $k[X]$ is integral over $A$.  Hence, if $X$ is normal, then $k[X]$ is the normalization of $A$.
\vskip 3pt
\emph{Hint:} Since $\psi$ is finite, every element $\overline{f}$ of $k[X]/I$ satisfies an integral expression over $k[\bA^m]/J \cong A/(\ker \beta)$ (even though these rings are not domains),
\[
\overline{f}^n + \overline{a}_{n-1} \overline{f}^{n-1} + \dots + \overline{a}_1 \overline{f}^{1} + \overline{a}_0 = 0 \in k[X]/I
\]
with $a_i \in A$.  %Consider the corresponding equation in $k[X]$ (without the $\overline{\bullet}$) and modify $a_0$ appropriately.
\vfill
Suppose $A \cong k[Y]$ for some affine variety $Y$.  We have an induced map of varieties:
\[
\xymatrix{
X \ar@{->>}[d] & \ar@{_{(}->}[l] Z_X(I) \ar@{->>}[d] \\
Y & \ar@{_{(}->}[l] Z(J).
}
\]
where the vertical surjections are because the ring maps are finite.
\newpage
\noindent
{\bf 5.}  Consider $X = \bA^1$ with $k[X] = k[x]$.  $I = (x(x-1))$.  $J = (0) \subseteq k = k[]$.  Compute $A$ and $Y$ and draw the diagram to show that $X \twoheadrightarrow Y$ is not bijective.
\vskip 130pt
\noindent
{\bf 6.}  Consider $X = \bA^1$ with $k[X] = k[x]$.  $I = (x^2)$.  $J = (0) \subseteq k$.  Compute $A$ and $Y$ and show that $X \to Y$ is bijective.
\vskip 130pt
\noindent
{\bf 7.}  Show that the map $X \to Y$ induces an isomorphism of quasi-affine varieties:
\[
X \setminus Z_X(I) \to Y \setminus \mathrm{Image}(Z(J)) = Y \setminus Z(\ker \beta).
\]
\emph{Hint:} Invert an element $f \in \ker \beta$ and show that $A[f^{-1}] \to k[X][f^{-1}]$ is an isomorphism.
\vskip 180pt
What we have shown so far is that $Y$ is obtained from $X$ excepted that we replaced $Z_X(I)$ (and possibly some tangent data around it as in {\bf 6.}) with $\mathrm{Image}(Z(J))$.  We do one more example.
\vskip 9pt
\noindent
{\bf 8.}  Consider $X = \bA^2$ with $k[X] = k[x,y]$ and set $I = (y)$.  Finally let $J = (0) \subseteq k[t]$ with the map $\psi : k[t] \to k[x,y]/(y) \cong k[x]$ sending $t \mapsto x^2$.  Compute $A$ and describe it geometrically in the case that $\mathrm{char} k \neq 2$.
\newpage
%\vskip 150pt
\noindent
{\bf 9.}  Suppose that $W$ is a non-normal affine variety with normalization (with respect to $k(W)$), and that $X \to W$ is the normalization (so that $k[W] \hookrightarrow k[X]$ is a finite map, both rings in $k(W)$).  Let $\mathfrak{c} = \{f \in k[W] \;|\; f \cdot k[X] \subseteq k[W] \}$.  Show that $\mathfrak{c}$ is an ideal in both $k[W]$ and $k[X]$.  It is called the \emph{conductor of $X \to W$ (or of $k[W] \subseteq k[X]$)}.
\vskip 200pt
\noindent
{\bf 10.}  If we set $I = \mathfrak{c} \subseteq k[X]$ and write $k[\bA^m]/J = k[W]/\mathfrak{c}$ with $\psi : k[W]/\mathfrak{c} \to k[X]/\mathfrak{c}$ the canonical map, show that the associated $A$ is isomorphic to $k[W]$.  In particular, \emph{every} non-normal affine variety is constructed by replacing a subvariety (and possibly some tangent data) of its normalization with another variety (with tangent data).
\vfill
\noindent
{\bf Extra.}  You might find it amusing to compute $A$ in the case of $X = \bA^2$, $I = (x)$ and $J = (0) \subseteq k$.
\end{document}

